\documentclass[pdflatex,sn-mathphys-num]{sn-jnl}

\usepackage{graphicx}%
\usepackage{multirow}%
\usepackage{amsmath,amssymb,amsfonts}%
\usepackage{amsthm}%
\usepackage{mathrsfs}%
\usepackage[title]{appendix}%
\usepackage{xcolor}%
\usepackage{textcomp}%
\usepackage{manyfoot}%
\usepackage{booktabs}%
\usepackage{algorithm}%
\usepackage{algorithmicx}%
\usepackage{algpseudocode}%
\usepackage{listings}%
\usepackage{lmodern}%
\usepackage{siunitx}
\usepackage{microtype}
\usepackage{derivative}
\usepackage[version=4]{mhchem}
\usepackage{cleveref}
\usepackage[spanish]{babel}

\renewcommand{\arraystretch}{1.2}
\sisetup{group-digits=true,
  group-separator={\,},
  separate-uncertainty
}

\theoremstyle{thmstyleone}%
\newtheorem{theorem}{Teorema}%
\newtheorem{proposition}[theorem]{Proposición}%

\theoremstyle{thmstyletwo}%
\newtheorem{example}{Ejemplo}%
\newtheorem{remark}{Observación}%

\theoremstyle{thmstylethree}%
\newtheorem{definition}{Definición}%

\DeclareSIUnit\angstrom{\text{\r{A}}}

\raggedbottom
\unnumbered% uncomment this for unnumbered level heads

\begin{document}

\title[]{}

\author{\fnm{Julian L.}
\sur{Avila-Martinez}}\email{jlavilam@udistrital.edu.co}

\affil{\orgdiv{Programa Académico de Física}, \orgname{Universidad Distrital
Francisco José de Caldas}, \orgaddress{\city{Bogotá D.C.},
\country{Colombia}}}

\abstract{
  Este artículo analiza las deficiencias epistemológicas y sociopolíticas en la
conceptualización de la nanotecnología. Se examina cómo la
historiografía del campo se fundamenta en mitos, ignorando que la verdadera
génesis disciplinar reside en los avances de la física fundamental.
Se demuestra que la métrica estándar de 1 a 100 nanómetros carece de rigor, al
eludir la fenomenología cuántica y termodinámica que define la escala
. Esta ceguera ontológica legitima un modelo de política pública
basado en el utilitarismo de la aplicación inmediata. Mediante el
caso de Colombia, se expone cómo la exigencia de innovación sin una inversión
proporcional en ciencia básica condena a las naciones en desarrollo a la
dependencia tecnológica y al estancamiento estructural.
}

\keywords{Nanotecnología, epistemología de la física,
política científica, fenomenología cuántica, ciencia básica, desarrollo
tecnológico.}

\maketitle

\section{Introducción}

El desarrollo de tecnologías de vanguardia exige una correspondencia estricta
entre el rigor analítico de la física teórica y la formulación de políticas de
Estado. Sin embargo, el campo de la nanotecnología exhibe una profunda fractura
entre su naturaleza física y su implementación institucional. Este
documento examina una cadena de errores categoriales que distorsionan el dominio
nanométrico, iniciando con la adopción de un mito historiográfico y culminando
en el colapso estructural de ecosistemas científicos en desarrollo. Se
argumenta que la imposición de métricas geométricas arbitrarias, en sustitución
de los criterios dictados por la mecánica cuántica, facilita una aproximación
puramente utilitaria de la ciencia. Al priorizar la aplicación
inmediata sobre el estudio de los principios fundamentales de la materia, las
políticas públicas generan un entorno de dependencia. El análisis se
centra en evidenciar cómo la suplantación del rigor físico por narrativas de
ingeniería macroscópica bloquea el progreso científico genuino, tomando como
paradigma empírico de este fracaso institucional el caso colombiano.


\section{Historia material}

La historiografía popular de la nanotecnología se fundamenta en un mito de
origen: la noción de que la disciplina surgió de la conferencia dictada por
Richard P. Feynman en 1959 \cite{feynmanTheresPlentyRoom2011}.
No obstante, el registro histórico revela que su influencia en la génesis del
campo fue nula. Aunque Feynman articuló experimentos mentales sobre la
miniaturización, su discurso carecía de los formalismos matemáticos y los
protocolos experimentales necesarios para guiar la manipulación atómica.

El análisis bibliométrico de Chris Toumey evidencia la desconexión entre
esta profecía y la práctica científica \cite{toumeyReadingFeynmanNanotechnology2008}.
La conferencia de 1959 generó apenas tres citas en la década de 1960 y cuatro en
la de 1970. El texto fue ignorado por la comunidad científica hasta ser
resucitado décadas después. Autores como K. Eric Drexler utilizaron a Feynman
como argumento de autoridad retrospectivo, buscando conferir el prestigio de un
laureado a agendas tecnológicas que carecían de rigor analítico.

La verdadera génesis de la nanotecnología reside en los avances formales de
la física de la materia condensada y la mecánica cuántica. La
manipulación atómica fue inaugurada por hitos instrumentales empíricos,
destacando la invención del microscopio de efecto túnel en 1981 por Gerd
Binnig y Heinrich Rohrer, y la posterior manipulación de átomos de
xenón por Don Eigler. Binnig y Rohrer no se fundamentaron en el discurso de
Feynman. Su desarrollo respondió estrictamente a la necesidad de resolver
anomalías en la física de superficies bajo vacío ultraalto.

De manera análoga, el descubrimiento de los fullerenos en 1985 por Kroto,
Smalley y Curl respondió a interrogantes fundamentales en la espectroscopía
molecular y la astroquímica. La base material de la disciplina se
construyó sobre el escrutinio de interacciones físicas. El progreso dependió
invariablemente del estudio de la naturaleza, demostrando que la innovación
tecnológica es consecuencia directa de la ciencia básica, no de premisas
teleológicas orientadas a la aplicación inmediata.


\section{La crisis ontológica y la arbitrariedad de la métrica nanométrica}

Más allá de su historiografía mitificada, la nanotecnología padece una profunda
crisis ontológica respecto a su definición formal. El estándar internacional, el
cual clasifica como nanotecnología a la manipulación de la materia con al menos
una dimensión entre 1 y 100 nanómetros, es fundamentalmente arbitrario y carece
de justificación física rigurosa. Si bien el límite inferior de 1
nanómetro se corresponde lógicamente con el diámetro cinético de los átomos
ligeros, la cota superior de 100 nanómetros constituye una frontera puramente
administrativa. Esta métrica estrictamente geométrica no
proporciona información alguna sobre la naturaleza física, la complejidad o el
comportamiento dinámico del sistema en cuestión.

Desde la perspectiva de la física fundamental, la transición del régimen
macroscópico, gobernado por la cinemática clásica, a la escala nanométrica,
dominada por la mecánica cuántica y las fluctuaciones térmicas, no ocurre en un
límite invariante de 100 nanómetros. Una definición físicamente
consistente exige como criterio la emergencia de fenómenos cuánticos o
estadísticos que diverjan del comportamiento macroscópico. Este
umbral de transición está dictado por parámetros intrínsecos del material,
tales como el radio de Bohr del excitón, el recorrido libre medio de los
electrones balísticos o la longitud de correlación de un estado fundamental
colectivo. Dependiendo del material y de la fenomenología observada,
estos efectos cuánticos pueden manifestarse a 2 nanómetros o extenderse por
encima de varios cientos de nanómetros.

Al subordinar la naturaleza cuántica de la materia a un umbral geométrico
simplista, la comunidad científica ha facilitado un blanqueamiento
terminológico. En este escenario, procesos químicos ordinarios son
reetiquetados como tecnología de vanguardia por el mero hecho de satisfacer
una restricción espacial.

Esta ausencia de rigor metodológico se exacerba a nivel estatal,
particularmente en ecosistemas científicos en desarrollo. En Colombia, el
intento de establecer un marco regulatorio derivó en una definición
operativa excesivamente laxa. El Consejo Nacional Asesor de Nanociencia y
Nanotecnología propuso clasificar como nanomaterial a cualquier sistema que
exhiba propiedades atribuibles a la escala hasta un límite de un micrómetro.
Adicionalmente, el criterio estipula que un material es
nanoparticulado si apenas el 10 \% de sus partículas presenta un tamaño igual
o inferior a 100 nanómetros, en franco contraste con el umbral del 50 \%
exigido por la Comisión Europea.

La adopción de una métrica que elude tanto el consenso global como las
fronteras estrictas de la física fundamental instituye un andamiaje
burocrático que diluye la integridad del campo. Esta laxitud conceptual
fomenta la apropiación discursiva, permitiendo que la ciencia de materiales
convencional se presente como nanotecnología avanzada sin demostrar un
verdadero dominio atómico.


\section{La omisión de la fenomenología cuántica y el puente aplicativo}

Para definir con rigor la nanotecnología, es imperativo abandonar las fronteras
geométricas arbitrarias y adoptar los principios físicos fundamentales que
gobiernan la materia a esta escala. La verdadera naturaleza del
dominio nanométrico radica en su fenomenología cuántica y en el predominio
absoluto de los efectos de superficie. Estos factores generan
comportamientos físicos no clásicos que divergen radicalmente de las
propiedades de la materia macroscópica. Al reducir los sistemas a
proporciones atómicas, los fenómenos de la mecánica cuántica y estadística
alteran drásticamente las propiedades electrónicas, ópticas y estructurales del
material. De forma análoga, la extrema proporción entre el área
superficial y el volumen dicta que las intensas fluctuaciones térmicas y las
fuerzas intermoleculares anulen la validez de la cinemática clásica.
Una definición rigurosa debe fundamentarse en la emergencia de estos fenómenos
físicos, regidos por parámetros intrínsecos como el radio de Bohr del excitón o
la longitud de correlación de un estado fundamental colectivo. Emplear
un límite estático de 100 nanómetros carece de todo significado físico.

Esta desvinculación de los fundamentos teóricos trasciende el simple error
académico para constituir la raíz de una profunda disfunción sistémica en la
estructuración del desarrollo tecnológico. Al despojar al campo de sus
realidades cuánticas y termodinámicas, reduciéndolo a un mero umbral
dimensional, se induce a la clase legislativa e industrial a percibir la
nanotecnología como una simple aplicación de ingeniería, ignorando su inherente
complejidad científica. Esta ceguera epistemológica alimenta el fracaso
definitivo de las políticas públicas. Existe un afán agresivo por implementar
aplicaciones tecnológicas avanzadas mientras se asfixia financieramente a la
ciencia teórica y experimental que las fundamenta. Al tratar a la
disciplina puramente como un emprendimiento aplicativo, los ecosistemas
científicos en desarrollo malinterpretan su naturaleza constitutiva.
Esta aproximación garantiza un estado de estancamiento estructural y perpetúa la
dependencia tecnológica de la nación.


\section{La falacia de la aplicación como principio y el colapso estructural}

La dependencia de la aplicación como principio rector constituye una ceguera
epistemológica profunda que suprime el verdadero avance científico. Al exigir
innovación tecnológica inmediata sin una inversión sustancial en la física
fundamental subyacente, la política pública atrapa al ecosistema científico en
un valle de la muerte, impidiendo la transición del conocimiento básico hacia
aplicaciones viables. Este enfoque inhibe la emergencia de
descubrimientos genuinos. Esto mantiene estática la producción científica
nacional, incluso cuando las directrices de financiación persiguen los
neologismos tecnológicos de turno.

En naciones en desarrollo, esta priorización utilitaria genera una peligrosa
ilusión de progreso. Existe la premisa falaz de que un país puede eludir el
trabajo riguroso de la ciencia teórica y adoptar instantáneamente disciplinas
avanzadas, como la nanotecnología, bajo un modelo de importación directa. En
realidad, sin una infraestructura robusta en física teórica, fisicoquímica y
educación científica estricta, la innovación es inalcanzable. Estas naciones
quedan relegadas a la posición de consumidores e importadores perpetuos de
tecnologías de caja negra desarrolladas en el Norte Global, imposibilitando su
consolidación como investigadores independientes o líderes científicos.

Colombia constituye un paradigma evidente de este fracaso sistémico. Las
políticas estatales exigen continuamente resultados comerciales altamente
aplicados. Se impulsa agresivamente el desarrollo de tecnologías convergentes y
de la Industria 4.0, mientras que la inversión nacional en investigación y
desarrollo permanece en niveles críticamente bajos, oscilando históricamente
entre el 0.20 \% y el 0.29 \% de su producto interno bruto. Esta
severa desfinanciación expone una profunda hipocresía estructural. Se percibe
un intenso apetito por el prestigio tecnológico inmediato, acoplado a un
desinterés absoluto por la educación fundamental y la investigación básica. El
marco institucional presiona continuamente a los investigadores para obtener
retornos económicos a corto plazo. Esto demuestra una incapacidad para
sacrificar recursos presentes en favor del bienestar científico a largo plazo.
En última instancia, la pretensión de cosechar los frutos de la
tecnología avanzada mientras se asfixian sus raíces teóricas garantiza que la
nación permanezca tecnológicamente colonizada y económicamente estancada.


\section{Conclusión}

La integridad de la nanotecnología como disciplina científica es incompatible
con las definiciones geométricas arbitrarias y el utilitarismo político.
La consolidación de un marco teórico riguroso requiere reconocer que
la transición hacia el dominio nanométrico está dictada exclusivamente por la
emergencia de fenómenos cuánticos y fluctuaciones termodinámicas, no por
fronteras espaciales estáticas. La ignorancia sistemática de estos
principios físicos ha legitimado políticas estatales que persiguen la aplicación
tecnológica inmediata mientras desmantelan la base analítica que la sustenta.
El caso de Colombia ilustra las consecuencias catastróficas de este
modelo, donde la falta de inversión en ciencia básica imposibilita la verdadera
innovación. Pretender dominar la materia a escala atómica eludiendo el
rigor teórico constituye un absurdo metodológico e institucional. La
verdadera autonomía tecnológica no se alcanza mediante la asimilación de un
discurso de innovación, sino a través de una inversión irrestricta en el estudio
profundo de las leyes fundamentales del universo. Sin este compromiso,
el progreso científico se reduce a una ilusión discursiva y administrativa.


\end{document}
